\documentclass[12pt,a4paper]{article}

\usepackage{fontspec}
\setmainfont{Tinos}
\setsansfont{Arimo}
\setmonofont{DejaVu Sans Mono}
\usepackage[russian]{babel}
\usepackage{geometry}
\geometry{left=38mm,right=25mm,top=38mm,bottom=25mm}
\usepackage{setspace}
\onehalfspacing
\usepackage{indentfirst}
\setlength{\parindent}{1.25cm}
\setlength{\emergencystretch}{3em}
\usepackage{xcolor}
\usepackage{listings}
\usepackage{booktabs}
\usepackage{array}
\usepackage{graphicx}
\usepackage{float}
\usepackage{caption}
\usepackage{titlesec}
\usepackage{needspace}
\usepackage{tikz}
\usetikzlibrary{arrows.meta,positioning,calc}
\usepackage{pgfplots}
\pgfplotsset{compat=1.18}
\usepackage{xurl}
\usepackage{hyperref}
\hypersetup{colorlinks=true,linkcolor=black,urlcolor=blue}
\usepackage{fancyhdr}
\pagestyle{fancy}
\fancyhf{}
\fancyfoot[C]{\thepage}
\renewcommand{\headrulewidth}{0pt}
\newcommand{\studentName}{Пантюхин Никита Олегович}
\newcommand{\studentGroup}{М8О-201БВ-25}
\newcommand{\checkerName}{Миронов Евгений Сергеевич}

\definecolor{codebg}{RGB}{247,247,247}
\definecolor{codeframe}{RGB}{205,205,205}
\definecolor{codecomment}{RGB}{70,110,70}
\definecolor{codekeyword}{RGB}{30,70,150}
\definecolor{headingone}{HTML}{1F3864}
\definecolor{headingtwo}{HTML}{2E5FA3}
\renewcommand{\lstlistingname}{Листинг}
\DeclareCaptionLabelSeparator{reportdash}{~---~}
\captionsetup[lstlisting]{labelsep=reportdash}
\titleformat{\section}
  {\color{headingone}\bfseries\fontsize{14}{17}\selectfont}
  {\thesection.}{0.5em}{}
\titleformat{\subsection}
  {\color{headingtwo}\bfseries\fontsize{13}{16}\selectfont}
  {\thesubsection.}{0.5em}{}
\titlespacing*{\section}{0pt}{18pt}{9pt}
\titlespacing*{\subsection}{0pt}{12pt}{6pt}
\lstdefinestyle{cstyle}{
  language=C,basicstyle=\ttfamily\footnotesize,
  keywordstyle=\color{codekeyword}\bfseries,
  commentstyle=\color{codecomment},backgroundcolor=\color{codebg},
  frame=single,rulecolor=\color{codeframe},breaklines=true,
  breakatwhitespace=false,showstringspaces=false,tabsize=4,
  columns=fullflexible,keepspaces=true,captionpos=b,
  aboveskip=0.7em,belowskip=0.9em
}
\lstdefinestyle{terminal}{
  basicstyle=\ttfamily\footnotesize,backgroundcolor=\color{codebg},
  frame=single,rulecolor=\color{codeframe},breaklines=true,
  breakatwhitespace=false,showstringspaces=false,
  columns=fullflexible,keepspaces=true,captionpos=b,
  aboveskip=0.7em,belowskip=0.9em
}

% Скриншоты автоматически подставляются при наличии файла в screenshots/.
% До добавления изображения на его месте отображается подписанное поле.
\newcommand{\reportscreenshot}[2]{%
  \begin{figure}[H]
  \centering
  \IfFileExists{screenshots/#1}{%
    \includegraphics[width=0.94\linewidth,height=7.0cm,keepaspectratio]{screenshots/#1}%
  }{%
    \fbox{\begin{minipage}[c][4.3cm][c]{0.86\linewidth}
      \centering\small\color{gray}
      Место для скриншота\\[0.6em]
      \texttt{screenshots/#1}
    \end{minipage}}%
  }
  \caption*{Рисунок~--- #2}
  \end{figure}
}

\begin{document}
\begin{titlepage}
\thispagestyle{empty}
\begin{center}
{\fontsize{10}{12}\selectfont\bfseries
МИНИСТЕРСТВО НАУКИ И ВЫСШЕГО ОБРАЗОВАНИЯ\\
РОССИЙСКОЙ ФЕДЕРАЦИИ}\\[0.2cm]
{\fontsize{10}{12}\selectfont
Федеральное государственное автономное образовательное учреждение высшего образования\\
\textbf{«Московский авиационный институт\\
(национальный исследовательский университет)»}}

\vspace{2.2cm}
{\fontsize{16}{19}\selectfont\bfseries ОТЧЁТ ПО ЛАБОРАТОРНОЙ РАБОТЕ}\\[0.15cm]
{\fontsize{16}{19}\selectfont\bfseries № 3}

\vspace{0.9cm}
\textbf{Дисциплина:}\\[0.15cm]
Операционные системы и архитектура компьютера

\vspace{0.7cm}
\textbf{Тема:}\\[0.15cm]
Многопоточная битоническая сортировка массива целых чисел

\vspace{2.0cm}
\begin{flushright}
\textbf{Выполнил:}\\[0.15cm]
студент группы \studentGroup\\[0.1cm]
\studentName\\[0.6cm]
\textbf{Проверил:}\\[0.15cm]
Преподаватель\\[0.1cm]
\checkerName
\end{flushright}
\vfill
Москва, 2026 г.
\end{center}
\end{titlepage}

\setcounter{page}{2}
\tableofcontents
\newpage

\section{Цель работы}
Цель работы --- получить практические навыки создания и синхронизации потоков в Linux. Необходимо реализовать на языке C многопоточную битоническую сортировку массива целых чисел, предусмотреть ограничение максимального количества потоков ключом запуска и продемонстрировать их фактическое количество средствами операционной системы. Также требуется исследовать зависимость ускорения и эффективности алгоритма от количества потоков, размера и типа входных данных.

\section{Краткое описание архитектуры решения}
Программа реализована на языке C с использованием библиотеки POSIX Threads (\texttt{pthread}). Проект собирается с помощью CMake. Вся реализация находится в одном исходном файле \texttt{main.c}: программа обрабатывает аргументы командной строки, создаёт массив, запускает рабочие потоки, измеряет время сортировки и проверяет порядок элементов.

Параметр \texttt{-t} задаёт максимальное количество одновременно участвующих потоков, включая главный. Параметр \texttt{-n} задаёт количество элементов, а необязательный параметр \texttt{-d} --- тип входных данных: \texttt{random}, \texttt{sorted}, \texttt{reverse} или \texttt{repeated}. Размер массива ограничен степенями двойки, начиная с двух элементов. Число используемых потоков не превышает $n/2$, поскольку на каждом шаге сети выполняется $n/2$ независимых сравнений.

Битоническая сортировка последовательно строит группы элементов размера \texttt{size} и сравнивает пары на расстоянии \texttt{distance}. На каждом этапе потоки получают непересекающиеся диапазоны пар, поэтому не изменяют одни и те же элементы одновременно. После этапа все потоки встречаются на барьере \texttt{pthread\_barrier\_wait()}, чтобы следующий этап начался только после завершения текущего. Массив находится в общей памяти процесса, отдельные процессы не создаются.

\clearpage
\subsection{Схема взаимодействия компонентов}
На рисунке показано распределение работы: главный поток также участвует в сортировке, остальные потоки создаются функцией \texttt{pthread\_create()}.

\begin{figure}[H]
\centering
\begin{tikzpicture}[
  block/.style={draw,rounded corners,minimum height=0.82cm,align=center,font=\small},
  line/.style={-{Latex[length=2.4mm]},thick},
  txt/.style={font=\scriptsize,fill=white,inner sep=1pt}
]
\node[block,minimum width=8.2cm] (main) at (0,0) {\texttt{main.c}: генерация массива, параметры \texttt{-t}, \texttt{-n}, \texttt{-d}};
\node[block,minimum width=8.2cm] (sort) at (0,-1.35) {Подготовка потоков и распределение пар сравнений};
\node[block,minimum width=3.55cm] (th0) at (-2.25,-2.9) {Главный поток\\\texttt{worker(tasks[0])}};
\node[block,minimum width=3.55cm] (thx) at (2.25,-2.9) {Дополнительные потоки\\\texttt{pthread\_create()}};
\node[block,minimum width=8.2cm] (barrier) at (0,-4.45) {\texttt{pthread\_barrier\_wait()}: ожидание всех потоков после этапа};
\node[block,minimum width=8.2cm] (finish) at (0,-5.8) {\texttt{pthread\_join()}, проверка результата и вывод времени};
\draw[line] (main) -- (sort);
\draw[line] (sort.south) |- (th0.north);
\draw[line] (sort.south) |- (thx.north);
\draw[line] (th0.south) |- (barrier.north);
\draw[line] (thx.south) |- (barrier.north);
\draw[line] (barrier) -- (finish);
\end{tikzpicture}
\caption*{Рисунок~--- Схема работы многопоточной сортировки}
\end{figure}

На каждом этапе выполняются независимые сравнения, а барьер обеспечивает правильный порядок этапов. Внутри этапа синхронизация между отдельными сравнениями не требуется.

\Needspace{0.55\textheight}
\section{Листинг ключевых фрагментов кода}
Основная работа выполняется функцией \texttt{worker()}. Внешний цикл увеличивает размер группы вдвое, внутренний цикл уменьшает расстояние между сравниваемыми элементами вдвое. Направление сравнения чередуется между группами.

\begin{lstlisting}[style=cstyle,caption={Сравнение элементов в рабочем потоке}]
void *worker(void *arg) {
    ThreadArgs *task = arg;

    for (long long size = 2; size <= task->n; size *= 2) {
        for (int distance = (int)(size / 2);
             distance > 0; distance /= 2) {
            for (int k = task->start; k < task->end; k++) {
                int i = (k / distance) * (2 * distance)
                        + (k % distance);
                int j = i + distance;
                int ascending = (i / size) % 2 == 0;

                if ((ascending && task->arr[i] > task->arr[j]) ||
                    (!ascending && task->arr[i] < task->arr[j])) {
                    int temp = task->arr[i];
                    task->arr[i] = task->arr[j];
                    task->arr[j] = temp;
                }
            }
            pthread_barrier_wait(task->barrier);
        }
    }
    return NULL;
}
\end{lstlisting}

\Needspace{0.32\textheight}
Каждому потоку передаются указатель на общий массив, размер массива, границы диапазона пар и указатель на барьер. Диапазоны формируются так, чтобы все $n/2$ пар были распределены между потоками без пересечений.

\begin{lstlisting}[style=cstyle,caption={Распределение диапазонов между потоками}]
for (int t = 0; t < thread_count; t++) {
    tasks[t].arr = arr;
    tasks[t].n = n;
    tasks[t].start = (int)((size_t)t * comparisons / thread_count);
    tasks[t].end = (int)((size_t)(t + 1) * comparisons / thread_count);
    tasks[t].barrier = &barrier;
}
\end{lstlisting}

Дополнительные потоки создаются с помощью \texttt{pthread\_create()}. Главный поток выполняет ту же функцию \texttt{worker()}, затем ожидает остальные потоки через \texttt{pthread\_join()}.

\begin{lstlisting}[style=cstyle,caption={Создание потоков и ожидание их завершения}]
for (int t = 1; t < thread_count; t++) {
    if (pthread_create(&threads[t - 1], NULL,
                       worker, &tasks[t]) != 0) {
        fprintf(stderr, "Thread creation failed\n");
        exit(EXIT_FAILURE);
    }
}

worker(&tasks[0]);

for (int t = 0; t < thread_count - 1; t++) {
    pthread_join(threads[t], NULL);
}
\end{lstlisting}

\Needspace{0.31\textheight}
Время измеряется монотонными часами \texttt{CLOCK\_MONOTONIC}: начало фиксируется перед созданием дополнительных потоков, а конец --- после их завершения. Генерация массива, проверка результата и вывод не входят в измеренный промежуток. При этом учитываются создание, синхронизация и завершение потоков.

\begin{lstlisting}[style=cstyle,caption={Начало измерения времени}]
struct timespec start_time, end_time;
clock_gettime(CLOCK_MONOTONIC, &start_time);
\end{lstlisting}

После выполнения функции \texttt{worker()} главным потоком и ожидания остальных потоков с помощью \texttt{pthread\_join()} фиксируется конечное время и рассчитывается разность.

\begin{lstlisting}[style=cstyle,caption={Окончание измерения времени}]
clock_gettime(CLOCK_MONOTONIC, &end_time);

double elapsed = (end_time.tv_sec - start_time.tv_sec) +
                 (end_time.tv_nsec - start_time.tv_nsec) / 1e9;
\end{lstlisting}

После сортировки программа последовательно проверяет, что каждый элемент не меньше предыдущего, и выводит \texttt{Sorted: yes} при выполнении этого условия.

\begin{lstlisting}[style=cstyle,caption={Проверка порядка элементов после сортировки}]
int sorted = 1;
for (int i = 1; i < n; i++) {
    if (arr[i - 1] > arr[i]) {
        sorted = 0;
        break;
    }
}
\end{lstlisting}

Полный исходный код программы находится в файле \texttt{lab3/src/main.c}.

\Needspace{0.36\textheight}
\section{Результаты работы}
Программа была собрана и проверена в среде WSL (Linux). Для исследования использовались массивы целых чисел размерами от 65\,536 до 4\,194\,304 элементов, а также 1, 2, 4 и 6 потоков. При создании случайных массивов использовалось фиксированное начальное значение генератора \texttt{srand(12345)}, что обеспечивает одинаковый набор данных при повторении запуска с тем же типом и размером массива.

\subsection{Консольный вывод и проверка потоков}
Ниже приведён фактический результат одного из запусков программы с шестью потоками. Время приведено по выводу терминала.

\begin{lstlisting}[style=terminal,caption={Пример успешной сортировки массива}]
$ ./build/lab3 -t 6 -n 1048576 -d repeated
Array size: 1048576
Requested threads: 6
Used threads: 6
Data type: repeated
Sorted: yes
Time: 0.294284 s
\end{lstlisting}

\Needspace{0.28\textheight}
Фактическое количество потоков было продемонстрировано стандартной командой Linux \texttt{ps -L}. В момент работы программы её вывод содержал шесть строк с одним PID и шестью разными TID.

\begin{lstlisting}[style=terminal,caption={Потоки процесса по данным команды ps -L}]
    PID     TID STAT COMMAND
  75244   75244 Rl+  lab3
  75244   75254 Rl+  lab3
  75244   75255 Rl+  lab3
  75244   75256 Rl+  lab3
  75244   75257 Rl+  lab3
  75244   75258 Rl+  lab3
\end{lstlisting}

Одинаковый PID подтверждает принадлежность потоков одному процессу, а разные значения TID --- существование отдельных потоков. Таким образом, при ключе \texttt{-t 6} фактически использовались шесть потоков, включая главный.

\subsection{Ускорение и эффективность}
Ускорение $S(p)$ определялось как отношение времени однопоточной сортировки к времени работы $p$ потоков. Эффективность $E(p)$ показывает, насколько полно используется доступный параллелизм:
\[
S(p)=\frac{T(1)}{T(p)},\qquad E(p)=\frac{S(p)}{p}\cdot100\%.
\]
Для массива из 1\,048\,576 случайных элементов каждую конфигурацию запускали пять раз. В таблице приведены все фактические времена в секундах.

\begin{table}[H]
\centering\small
\renewcommand{\arraystretch}{1.22}
\begin{tabular}{c r r r r}
\toprule
\textbf{Запуск} & \textbf{1 поток} & \textbf{2 потока} & \textbf{4 потока} & \textbf{6 потоков} \\
\midrule
1 & 0,994904 & 0,474835 & 0,250886 & 0,336671 \\
2 & 0,972518 & 0,498422 & 0,267290 & 0,313168 \\
3 & 0,794891 & 0,426889 & 0,289247 & 0,330345 \\
4 & 0,837853 & 0,548274 & 0,287952 & 0,320280 \\
5 & 0,997526 & 0,537042 & 0,386954 & 0,327392 \\
\bottomrule
\end{tabular}
\caption*{Таблица~--- Результаты пяти повторных измерений, с}
\end{table}

Чтобы уменьшить влияние разовых колебаний, для каждого количества потоков выбрана медиана пяти измерений. Для ускорения используется отношение медианного времени одного потока к медианному времени соответствующего числа потоков.

\begin{table}[H]
\centering\small
\renewcommand{\arraystretch}{1.22}
\begin{tabular}{c r r r}
\toprule
\textbf{Потоки} & \textbf{Медиана, с} & \textbf{Ускорение} & \textbf{Эффективность, \%} \\
\midrule
1 & 0,972518 & 1,00 & 100,00 \\
2 & 0,498422 & 1,95 & 97,56 \\
4 & 0,287952 & 3,38 & 84,43 \\
6 & 0,327392 & 2,97 & 49,51 \\
\bottomrule
\end{tabular}
\caption*{Таблица~--- Ускорение и эффективность по медианным временам}
\end{table}

\begin{figure}[H]
\centering
\begin{tikzpicture}
\begin{axis}[
 width=0.94\linewidth,height=6.0cm,
 xlabel={Количество потоков},ylabel={Ускорение, раз},
 xmin=0.8,xmax=6.2,ymin=0,ymax=6.5,
 xtick={1,2,4,6},grid=major,grid style={gray!20},
 legend pos=north west,
 legend style={font=\footnotesize,draw=none,fill=none},
 tick label style={font=\small},label style={font=\small}]
\addplot[color=headingtwo,mark=*,thick] coordinates {(1,1) (2,1.9512) (4,3.3774) (6,2.9705)};
\addlegendentry{Измеренное ускорение}
\addplot[color=gray,dashed,thick] coordinates {(1,1) (2,2) (4,4) (6,6)};
\addlegendentry{Идеальное ускорение}
\end{axis}
\end{tikzpicture}
\caption*{Рисунок~--- Зависимость ускорения от количества потоков}
\end{figure}

\begin{figure}[H]
\centering
\begin{tikzpicture}
\begin{axis}[
 width=0.94\linewidth,height=5.9cm,
 xlabel={Количество потоков},ylabel={Эффективность, \%},
 xmin=0.8,xmax=6.2,ymin=0,ymax=110,
 xtick={1,2,4,6},ytick={0,20,40,60,80,100},
 grid=major,grid style={gray!20},
 tick label style={font=\small},label style={font=\small}]
\addplot[color=headingtwo,mark=square*,thick] coordinates {(1,100) (2,97.5597) (4,84.4340) (6,49.5083)};
\end{axis}
\end{tikzpicture}
\caption*{Рисунок~--- Зависимость эффективности от количества потоков}
\end{figure}

На данном размере массива наименьшее медианное время получено при четырёх потоках. Шесть потоков работали несколько медленнее: накладные расходы на создание и синхронизацию потоков, а также особенности выполнения в ОС могут снижать дополнительный выигрыш от параллелизма.

\subsection{Влияние размера входного массива}
Отдельно проведено по одному измерению для четырёх размеров случайного массива при каждом количестве потоков. Результаты приведены в секундах.

\begin{table}[H]
\centering\small
\renewcommand{\arraystretch}{1.22}
\begin{tabular}{r r r r r}
\toprule
\textbf{Элементов} & \textbf{1 поток} & \textbf{2 потока} & \textbf{4 потока} & \textbf{6 потоков} \\
\midrule
65\,536 & 0,042258 & 0,029752 & 0,024685 & 0,024834 \\
262\,144 & 0,196706 & 0,125991 & 0,089934 & 0,083893 \\
1\,048\,576 & 0,973696 & 0,475099 & 0,284199 & 0,339342 \\
4\,194\,304 & 4,903357 & 1,966317 & 1,108901 & 1,102667 \\
\bottomrule
\end{tabular}
\caption*{Таблица~--- Время сортировки в зависимости от размера массива, с}
\end{table}

\begin{figure}[H]
\centering
\begin{tikzpicture}
\begin{axis}[
 width=0.94\linewidth,height=6.6cm,
 xlabel={Размер массива, элементов},ylabel={Время, с},
 symbolic x coords={65536,262144,1048576,4194304},
 xtick=data,xticklabels={65 тыс.,262 тыс.,1 млн,4,2 млн},
 ymin=0,grid=major,grid style={gray!20},
 legend style={at={(0.02,0.98)},anchor=north west,font=\scriptsize,draw=none,fill=none},
 tick label style={font=\small},label style={font=\small}]
\addplot[mark=*,thick,color=black] coordinates {(65536,0.042258) (262144,0.196706) (1048576,0.973696) (4194304,4.903357)};
\addlegendentry{1 поток}
\addplot[mark=square*,thick,color=headingtwo] coordinates {(65536,0.029752) (262144,0.125991) (1048576,0.475099) (4194304,1.966317)};
\addlegendentry{2 потока}
\addplot[mark=triangle*,thick,color=green!50!black] coordinates {(65536,0.024685) (262144,0.089934) (1048576,0.284199) (4194304,1.108901)};
\addlegendentry{4 потока}
\addplot[mark=diamond*,thick,color=orange!80!black] coordinates {(65536,0.024834) (262144,0.083893) (1048576,0.339342) (4194304,1.102667)};
\addlegendentry{6 потоков}
\end{axis}
\end{tikzpicture}
\caption*{Рисунок~--- Зависимость времени сортировки от размера массива}
\end{figure}

В таблице ниже для каждого размера представлены ускорение и эффективность шести потоков относительно одного, рассчитанные по соответствующим одиночным измерениям.

\begin{table}[H]
\centering\small
\renewcommand{\arraystretch}{1.22}
\begin{tabular}{r r r}
\toprule
\textbf{Элементов} & \textbf{Ускорение, 6 потоков} & \textbf{Эффективность, \%} \\
\midrule
65\,536 & 1,70 & 28,36 \\
262\,144 & 2,34 & 39,08 \\
1\,048\,576 & 2,87 & 47,82 \\
4\,194\,304 & 4,45 & 74,11 \\
\bottomrule
\end{tabular}
\caption*{Таблица~--- Ускорение и эффективность для различных размеров массивов}
\end{table}

При увеличении массива с 65\,536 до 4\,194\,304 элементов преимущество параллельной обработки становится заметнее. На наибольшем массиве один поток выполнял сортировку за 4,903357 с, а шесть потоков --- за 1,102667 с, то есть в данном запуске ускорение составило около 4,45 раза. Для небольшого массива затраты на работу с потоками составляют более значимую долю общего времени.

\subsection{Влияние типа входных данных}
Для массива из 1\,048\,576 элементов дополнительно сравнивались предварительно отсортированные, обратно отсортированные и содержащие повторения данные. Для каждой конфигурации выполнено по одному запуску.

\begin{table}[H]
\centering\small
\renewcommand{\arraystretch}{1.22}
\begin{tabular}{l r r r r}
\toprule
\textbf{Тип массива} & \textbf{1 поток} & \textbf{2 потока} & \textbf{4 потока} & \textbf{6 потоков} \\
\midrule
\texttt{sorted} & 0,996246 & 0,405243 & 0,380148 & 0,312504 \\
\texttt{reverse} & 0,709495 & 0,418982 & 0,444729 & 0,467639 \\
\texttt{repeated} & 1,324132 & 1,406348 & 0,852322 & 0,624758 \\
\bottomrule
\end{tabular}
\caption*{Таблица~--- Время сортировки различных входных данных, с}
\end{table}

Для оценки влияния типа входных данных на масштабирование также рассчитаны ускорение и эффективность шести потоков относительно одного. Эти оценки основаны на одиночных измерениях.

\begin{table}[H]
\centering\small
\renewcommand{\arraystretch}{1.22}
\begin{tabular}{l r r}
\toprule
\textbf{Тип массива} & \textbf{Ускорение, 6 потоков} & \textbf{Эффективность, \%} \\
\midrule
\texttt{sorted} & 3,19 & 53,13 \\
\texttt{reverse} & 1,52 & 25,29 \\
\texttt{repeated} & 2,12 & 35,32 \\
\bottomrule
\end{tabular}
\caption*{Таблица~--- Масштабирование сортировки для разных типов данных}
\end{table}

В битонической сортировке последовательность этапов и число сравнений определяются размером массива, а не исходным порядком чисел. Различия наблюдаемого времени могут зависеть от числа обменов, поведения процессора, синхронизации и текущей нагрузки системы. Так как эти конфигурации измерялись однократно, результаты показывают наблюдения, но не доказывают устойчивое преимущество какого-либо типа массива или числа потоков.

\clearpage
\subsection{Скриншоты результатов}
В разделе предусмотрены места для фактических скриншотов терминала. Файлы \texttt{run.png}, \texttt{threads.png} и \texttt{measurements.png} необходимо положить в папку \texttt{report/screenshots/}; при следующей сборке они автоматически появятся на своих местах.

\reportscreenshot{run.png}{Успешный запуск многопоточной сортировки}
\reportscreenshot{threads.png}{Наблюдение шести потоков с помощью команды ps -L}
\reportscreenshot{measurements.png}{Измерения времени сортировки при разном количестве потоков}

\Needspace{0.24\textheight}

\section{Выводы}

После выполнения этой лабораторной я лучше понял, как устроена
многопоточность в Linux и чем потоки отличаются от отдельных процессов.
На практике я разобрался, как создавать потоки с помощью
\texttt{pthread\_create()}, ожидать их завершения через
\texttt{pthread\_join()} и синхронизировать работу с помощью барьеров.

При реализации битонической сортировки я научился распределять
вычисления между несколькими потоками и проверять их количество
средствами операционной системы. Также я увидел, что увеличение
числа потоков не всегда приводит к ускорению программы, поскольку
на производительность влияют размер массива и затраты на синхронизацию.

Благодаря этой работе я стал лучше понимать принципы параллельного
выполнения задач и способы оценки его эффективности.
\end{document}
